
\documentclass[11pt]{article}
\usepackage[margin=0.7in]{geometry}
\usepackage{booktabs}
\usepackage{graphicx}
\usepackage{float}
\usepackage{array}
\usepackage{hyperref}
\hypersetup{colorlinks=true, linkcolor=blue, urlcolor=blue}

\title{Model Comparison: Unconditional, LightGBM, MLP, and TP2-MV-ALiBi Transformer\\
\large All Models Under Identical Causal Selection and Metrics -- Freshly Run, Not Pulled From Prior Reports}
\author{SPY American Call TP2 Project -- ICAIF Submission Support}
\date{\today}

\begin{document}
\sloppy
\maketitle

\section{Executive Summary}
Every number in this report comes from a script run this session --
\texttt{27\_train\_lightgbm\_final.py}, \texttt{28\_train\_mlp\_final.py},
\texttt{23\_train\_mv\_alibi\_transformer\_v3\_final.py}, and
\texttt{24\_train\_mv\_alibi\_transformer\_v4\_final.py} -- not pulled from any earlier PDF. All four
learned models and the unconditional baseline are evaluated identically: isotonic calibration fit on
validation predictions only, causal $p\ge0.5$ selection (a fixed constant, zero tuning, no reference to any
other trade's prediction), and daily-grouped Sharpe, Sortino, and max drawdown -- the same formulas, same
code pattern, across every model. Sortino is reported for every model, not only the transformer.

\begin{table}[H]
\centering
\small
\begin{tabular}{lrrrrrrr}
\toprule
Model & Trades & Total P\&L & Win rate & Sharpe & Sortino & Max DD & Pos. years \\
\midrule
Unconditional (no ML) & 250,004 & -\$12,521,429 & 12.9\% & -1.01 & -2.12 & -\$24,505,562 & 2/6 \\
LightGBM (T1+K2) & 19,636 & \$682,411 & 50.6\% & 1.25 & 1.12 & -\$254,256 & 5/6 \\
MLP ($F^*$) & 14,258 & \$1,411,596 & 54.0\% & 2.60 & 6.89 & -\$215,118 & 6/6 \\
Transformer v3-final (1 seed) & 15,015 & \$570,128 & 61.0\% & 1.51 & 1.24 & -\$267,840 & 6/6 \\
Transformer v4-final (5-seed ens.) & 14,648 & \$703,705 & 60.1\% & 2.36 & 3.72 & -\$90,991 & 6/6 \\
\bottomrule
\end{tabular}
\caption{All models, 2021--2026 test years combined. LightGBM and MLP restrict to T1+K2 (the structural
short-D-leg prior); LightGBM's full 4-strategy (T1/T2/K1/K2) result is reported separately below as
evidence for that restriction.}
\end{table}

\textbf{Headline:} every learned model beats the unconditional floor by a wide margin (floor: Sharpe
-1.01, Sortino -2.12, total P\&L -\$12,521,429).
Under this identical, causal methodology, single-seed \textbf{MLP on $F^*$ has the highest Sharpe and
Sortino of any single model tested} (2.60 /
6.89), ahead of the 5-seed transformer ensemble
(2.36 / 3.72) and
LightGBM (1.25 / 1.12).
\textbf{This comparison is not apples-to-apples on ensembling, and correcting for that reverses the
ranking:} once both models are given the identical 5-seed ensembling treatment, the transformer ensemble
outperforms the MLP ensemble on Sharpe (2.36 vs.\ 1.66) -- see the ``Seed-Count Sensitivity'' section
below, which also shows the underlying single-seed gap is not statistically significant. Reported honestly
in both directions, not adjusted to favor either model.

\section{Every Model in This Comparison Is T1+K2 Only -- Why LightGBM's Table Looks Different}
\textbf{To be unambiguous: every number in Table 1 (LightGBM, MLP, both transformer versions) is T1+K2
only.} MLP and the transformer were never trained on, or exposed to, T2/K1 data at all -- their underlying
dataset (\texttt{tp2\_t1k2\_token\_dataset.parquet}) is T1+K2-restricted from the start. LightGBM is the one
model in this comparison that was \emph{additionally} trained on all four strategies (T1/T2/K1/K2), to
preserve its role as the project's feature-discovery model -- so it is also the only model for which a
broader, unrestricted result exists to show. That unrestricted number is included below purely as
justification for the T1+K2 structural prior (both strategies short the shorter-dated leg, decided
theoretically before looking at performance, not post-hoc) -- \textbf{it is not part of the headline
comparison in Table 1, and LightGBM's row in Table 1 is exactly as T1+K2-restricted as every other model's.}
The unrestricted result confirms the prior empirically:
\begin{table}[H]
\centering
\begin{tabular}{lrrrr}
\toprule
LightGBM scope & Trades & Total P\&L & Sharpe & Positive years \\
\midrule
All 4 strategies (T1/T2/K1/K2) & 124,467 & -\$37,232,222 & -1.76 & 1/6 \\
T1+K2 only & 19,636 & \$682,411 & 1.25 & 5/6 \\
\bottomrule
\end{tabular}
\end{table}
The unrestricted result loses heavily (T2 and K1 both short the longer-dated leg, giving more time for
losses to develop); T1+K2 alone is solidly profitable under the identical selection rule.

\textbf{Verified, not assumed:} because LightGBM's feature dataset technically starts from a larger
episode population (226,577 episodes, covering all four strategies) than MLP/the transformer's
T1+K2-only dataset (177,487 episodes), it was checked directly whether LightGBM's T1 and K2 sub-models
end up training and evaluating on the exact same episodes as the other two models, not just a
similarly-sized but different set. They do: LightGBM's T1 rows (177,487, after joining to valid T1 labels)
and MLP/the transformer's T1 rows are the identical episode-ID set (zero symmetric difference); same for
K2. Base profitability rates match exactly (T1: 0.2533 both sources; K2: 0.2367 both sources). Fold
train/validation/test year boundaries are also textually identical across all three training scripts.
LightGBM's T1 and K2 sub-models are trained as fully independent binary classifiers (separate model
instances, separate \texttt{fit()} calls) -- there is no parameter sharing or information leakage from the
T2/K1 sub-models into the T1/K2 ones. The remaining differences between LightGBM and MLP/the transformer
are architectural (per-strategy independent binary classifiers vs.\ a single model handling both
strategies) and feature-set (LightGBM uses the full $\sim$136-feature candidate set, consistent with its
role as the feature-discovery model; MLP and the transformer use the post-selection $F^*$, 48 features) --
both deliberate, disclosed choices, not scope inconsistencies.

\section{Year-by-Year P\&L / Sharpe / Sortino}
Format per cell: P\&L / Sharpe / Sortino.
\begin{table}[H]
\centering
\scriptsize
\begin{tabular}{lrrrrr}
\toprule
Year & Unconditional & LightGBM (T1+K2) & MLP & Transformer v3-final & Transformer v4-final \\
\midrule
2021 & \$3,235,718 & \$168,555/0.66/0.70 & \$87,653/0.63/1.10 & \$185,503/1.16/0.91 & \$161,371/1.55/3.23 \\
2022 & -\$2,583,859 & \$54,095/4.53/37.55 & \$6,884/6.59/35.51 & \$26,379/1.06/1.06 & \$28,821/1.69/1.78 \\
2023 & -\$1,676,247 & \$381,339/6.44/77.22 & \$764,933/4.02/19.87 & \$76,045/4.21/5.70 & \$123,772/5.93/15.63 \\
2024 & -\$5,932,142 & -\$37,474/-0.63/-0.51 & \$73,126/1.90/1.89 & \$72,525/2.02/1.82 & \$28,164/0.66/0.52 \\
2025 & -\$10,495,227 & \$38,622/0.38/0.26 & \$448,329/5.11/89.59 & \$205,897/3.95/6.53 & \$226,194/4.15/6.14 \\
2026 & \$4,930,329 & \$77,276/4.32/51.63 & \$30,671/6.38/$\infty$ & \$3,779/2.20/6.54 & \$135,382/5.35/488.98 \\
\bottomrule
\end{tabular}
\end{table}

\begin{figure}[H]
\centering
\includegraphics[width=0.95\linewidth]{figures/model_comparison_yearly_pnl.pdf}
\caption{Year-by-year normalized P\&L, all four learned models, identical causal $p\ge0.5$ selection.}
\end{figure}

\section{Technical Audit: What Was Missing Per Model, and What Was Fixed}
Before this pass, the models were not on comparable technical footing. Each gap below was found by
inspecting the actual script, not assumed.

\begin{itemize}
\item \textbf{LightGBM (\texttt{ExpiryMLCorrection/02\_parallel\_binary\_lgbm.py}, original):} no
probability calibration (a \texttt{class\_weight="balanced"} model's raw \texttt{predict\_proba} output is
not a well-calibrated probability); selection via \texttt{tune\_threshold} (validation-P\&L-maximizing
search, the same mechanism already retired from the transformer for producing unstable results); Sharpe
computed from \emph{trade-level}, not daily-grouped, returns -- the original script's own comment admits
this ("approximate ... not daily grouped here"), which is a materially different and incorrect statistic;
no Sortino or max drawdown computed at all. \textbf{Fixed in \texttt{27\_train\_lightgbm\_final.py}:}
added isotonic calibration (validation-only fit), switched to causal $p\ge0.5$ selection, corrected Sharpe to
daily-grouped, added Sortino and max drawdown.
\item \textbf{MLP (\texttt{02\_train\_mlp\_baseline.py}, original):} no calibration; selection via the same
\texttt{tune\_threshold} mechanism; Sharpe was already correctly daily-grouped; no Sortino or max drawdown.
\textbf{Fixed in \texttt{28\_train\_mlp\_final.py}:} added isotonic calibration, switched to causal
$p\ge0.5$ selection, added Sortino and max drawdown.
\item \textbf{Transformer:} already fixed in the prior pass (\texttt{shipped\_final\_report.pdf}) --
included here for reference, not re-run again in this report.
\item \textbf{Unconditional baseline:} newly computed for this report (not previously tracked against the
same 2021--2026 test-year window with matching daily-grouped metrics) -- gives an honest floor for the
comparison table above.
\end{itemize}

\section{Seed-Count Sensitivity: Ensembling Helps the Transformer, Hurts MLP}
The headline MLP-vs-transformer comparison above compares single-seed MLP against a 5-seed-ensembled
transformer (v4-final) -- a confound raised directly by the professor (meeting transcript, 2026-07-15/16,
[12:23]): ``You can see a little bit of that because you are averaging models and only then you get better
performance.'' To make the comparison fair, both models were tested across seed counts, using the identical
ensembling methodology (raw probabilities averaged across seeds, then one isotonic calibration).
\begin{table}[H]
\centering
\small
\begin{tabular}{lrrrrrr}
\toprule
Model & Trades & Total P\&L & Sharpe & Sortino & Max DD & Pos. years \\
\midrule
Transformer, 1 seed & 15,015 & \$570,128 & 1.51 & 1.24 & -\$267,840 & 6/6 \\
Transformer, 5 seeds & 14,648 & \$703,705 & 2.36 & 3.72 & -\$90,991 & 6/6 \\
Transformer, 10 seeds & 12,772 & \$746,544 & 1.81 & 1.59 & -\$297,707 & 6/6 \\
MLP, 1 seed & 14,258 & \$1,411,596 & 2.60 & 6.89 & -\$215,118 & 6/6 \\
MLP, 5 seeds & 24,977 & \$3,433,936 & 1.66 & 6.12 & -\$1,207,038 & 4/6 \\
\bottomrule
\end{tabular}
\end{table}
\textbf{Ensembling moves the two models in opposite directions.} For the transformer, 1$\to$5 seeds
improves Sharpe (1.51$\to$2.36) and drawdown (-\$268K$\to$-\$91K); 5$\to$10 seeds partially reverses this
(Sharpe 2.36$\to$1.81, drawdown back out to -\$298K) -- 5 seeds is a sweet spot, not a monotonic
more-is-better relationship, though total P\&L keeps climbing at every step (\$570K$\to$\$704K$\to$\$747K).
For MLP, 1$\to$5 seeds \textbf{hurts} Sharpe (2.60$\to$1.66) and consistency (6/6$\to$4/6 positive years),
even as total P\&L triples (\$1.41M$\to$\$3.43M) -- win rate collapses from 54.0\% to 33.0\% as ensembling
lets roughly 75\% more, lower-quality trades cross the $p\ge0.5$ bar.

\textbf{Once the comparison is made fair -- both single-seed, or both ensembled -- ``MLP beats the
transformer'' does not hold up cleanly:} single-seed MLP beats single-seed transformer (Sharpe 2.60 vs.\
1.51), but the 5-seed-ensembled transformer beats the 5-seed-ensembled MLP (2.36 vs.\ 1.66). \textbf{The
original headline comparison in this report's Executive Summary was, in this specific sense, comparing the
wrong things to each other.}

\section{Discussion}
Why MLP wins single-seed-vs-single-seed, and why ensembling affects the two models so differently, remain
open questions -- plausible, unconfirmed explanations: (1) MLP's calibrated probabilities may simply cross
0.5 less often on weaker candidates in a single fit, producing a pickier trade set that ensembling then
dilutes (5 diverse MLP fits pull individual seeds' more confident, extreme predictions toward the mean,
letting more borderline trades through after recalibration); (2) the transformer's added complexity
(historical cross-attention over up to 64 prior episodes) may benefit specifically from averaging across
seeds in a way a much smaller, faster-converging model does not; (3) six test years is a small sample.
\textbf{Formal significance testing was run on the single-seed comparison (bootstrap CI + sign test,
identical methodology to this project's v2-vs-v3 test): neither MLP-vs-v3-final nor MLP-vs-v4-final is
statistically significant} (sign-test $p$=0.344 and $p$=0.656 respectively; both 95\% CIs on the mean P\&L
delta include zero; one year, 2023, drives most of the aggregate gap in both comparisons).
\textbf{Bottom line for the paper: neither ``MLP is better'' nor ``the transformer is better'' is a
defensible, unqualified claim from this evidence -- the honest framing is that the two are close, the
ranking depends on ensembling choices that were not controlled for in the original comparison, and six
years cannot yet distinguish them statistically.}

\section{Scope and Limitations}
\begin{itemize}
\item No other tree-based models (XGBoost, CatBoost, Random Forest) exist in this codebase to compare --
LightGBM is the only tree model used in this project.
\item Statistical significance (bootstrap CI, sign test) has not been computed for this specific 4-model
comparison; it exists elsewhere in this project only for the v2-vs-v3 transformer comparison.
\item Transaction costs are not modeled for any model here, consistent with every other result in this project.
\item This report does not modify \texttt{ExpiryMLCorrection/02\_parallel\_binary\_lgbm.py},
\texttt{02\_train\_mlp\_baseline.py}, or any transformer script -- all fixes are in new, separate files.
\end{itemize}

\end{document}
